\begin{enumerate}
  \item Анализ теоретических основ молекулярной спектроскопии высокого разрешения показывает, что
  эффективный гамильтониан, приближение Борна--Оппенгеймера и теория изотопозамещения образуют
  единую методическую основу для исследования изотопологов молекул различной симметрии. Общность
  потенциальной и дипольной поверхностей в декартовом представлении позволяет переносить информацию от
  хорошо изученной материнской молекулы к менее изученным изотопологам, а масс-зависимые изменения
  нормальных координат, моментов инерции и молекулярно-фиксированных осей определяют наблюдаемые
  изотопические сдвиги и изменения интенсивностей.

  \item Для молекулы GeH$_4$ использование аналитически симметризованных колебательных функций и
  матричных элементов эффективного гамильтониана позволило получить физически согласованный набор
  спектроскопических параметров. Для ${}^{76}$GeH$_4$ подгонка 31 экспериментального центра полос
  описывается 27 параметрами со среднеквадратичным отклонением $d_{\textrm{rms}}=0.13$ см$^{-1}$, а для
  изотопологов ${}^{74}$GeH$_4$, ${}^{73}$GeH$_4$, ${}^{72}$GeH$_4$ и ${}^{70}$GeH$_4$ совместная подгонка
  77 центров полос описывается 42 параметрами с $d_{\textrm{rms}}=0.18$ см$^{-1}$. Это подтверждает, что
  изотопические ограничения и учет локально-модового характера растягивающих колебаний обеспечивают
  устойчивое описание колебательной структуры пяти стабильных изотопологов германа.

  \item Для молекулы SiH$_4$ полуэмпирическое решение обратной спектроскопической задачи обеспечивает
  более точное описание известных центров полос, чем расчет, основанный только на силовом поле. Для
  ${}^{28}$SiH$_4$ полученный набор параметров воспроизводит исходные центры полос со
  среднеквадратичным отклонением $d_{\textrm{rms}}=0.28$ см$^{-1}$, тогда как альтернативный расчет дает
  $d_{\textrm{rms}}=2.21$ см$^{-1}$. Для ${}^{29}$SiH$_4$ и ${}^{30}$SiH$_4$ совместная подгонка 32
  экспериментальных центров полос с учетом изотопических соотношений приводит к отклонению
  $d_{\textrm{rms}}=0.06$ см$^{-1}$. Полученные параметры позволяют рассчитать 6375 центров
  колебательных полос, то есть по 2125 центров для каждого из трех стабильных изотопологов.

  \item В области фундаментальной полосы $\nu_2$ молекулы D$_2{}^{32}$S отнесение 1742 переходов с
  $J_{\max}=30$ и $K_{a,\max}=21$ позволило определить 535 энергий состояния $(010)$. Уточнение
  параметров основного состояния с использованием 270 комбинационных разностей уменьшает систематические
  отклонения при больших $J$ и $K_a$, а последующая подгонка 34 параметров эффективного гамильтониана
  Ватсона воспроизводит энергии состояния $(010)$ со среднеквадратичным отклонением
  $1.31\times10^{-4}$ см$^{-1}$. Такой результат существенно расширяет известную энергетическую
  информацию о фундаментальной деформационной полосе D$_2$S.

  \item Для D$_2{}^{32}$S установлено, что корректное определение абсолютных интенсивностей полосы
  $\nu_2$ требует учета изменения состава газовой смеси вследствие водород-дейтериевого обмена.
  Использование изотопических соотношений для главного параметра эффективного дипольного момента и
  анализ пробных линий дают оптимальную оценку парциального давления D$_2$S
  $P_{\textrm{part}}=2.378\pm0.036$ мбар. После введения этой поправки набор из восьми параметров
  эффективного оператора дипольного момента воспроизводит 280 экспериментальных интенсивностей со
  среднеквадратичным отклонением 2.5~\%, что сопоставимо с экспериментальной неопределенностью анализа
  линий.

  \item Для молекулы HD$^{32}$S в области полосы $\nu_2$ отнесение 2684 переходов с $J_{\max}=25$ и
  $K_{a,\max}=17$ позволило определить 415 энергий состояния $(010)$ и получить набор из 33 параметров
  эффективного гамильтониана со среднеквадратичным отклонением $1.1\times10^{-4}$ см$^{-1}$. Для
  молекулы HD$^{34}$S было отнесено 1212 переходов с $J_{\max}=21$ и $K_{a,\max}=14$, определено 263
  энергии состояния $(010)$ и получен набор параметров, воспроизводящий исходные энергии со
  среднеквадратичным отклонением $1.6\times10^{-4}$ см$^{-1}$. Следовательно, расширение набора
  отнесенных линий обеспечивает более полное и точное описание фундаментальной полосы $\nu_2$ двух
  монодейтерированных изотопологов сероводорода.

  \item Для HD$^{32}$S учет понижения симметрии при переходе H$_2$S$\rightarrow$HDS и поворота
  молекулярно-фиксированной системы координат позволил использовать теорию изотопозамещения для оценки
  параметров эффективного дипольного момента и парциального давления HDS в образце. Полученное значение
  $P_{\textrm{part}}=1.77\pm0.10$ мбар соответствует доле HD$^{32}$S около 39.4~\%, а набор параметров
  эффективного дипольного момента воспроизводит 418 экспериментальных интенсивностей со
  среднеквадратичным отклонением 4.7~\%. Для HD$^{34}$S ограниченный анализ наиболее сильных линий дает
  отклонение около 4.0~\%, что подтверждает применимость единой схемы анализа интенсивностей для
  изотопологов HDS.

  \item В результате работы сформированы согласованные наборы спектроскопических данных: центры
  колебательных полос изотопологов GeH$_4$ и SiH$_4$, параметры эффективного гамильтониана и
  эффективного дипольного момента для дейтерированных изотопологов сероводорода, а также списки
  переходов полосы $\nu_2$ D$_2{}^{32}$S, HD$^{32}$S и HD$^{34}$S с рассчитанными индивидуальными
  интенсивностями. Эти данные могут использоваться при моделировании спектров высокого разрешения,
  уточнении потенциальных и дипольных поверхностей и подготовке расчетных материалов для
  спектроскопических баз данных.
\end{enumerate}
